% !TeX root = ../main.tex
\chapter*{DANH MỤC TỪ VIẾT TẮT VÀ THUẬT NGỮ}
\addcontentsline{toc}{chapter}{DANH MỤC TỪ VIẾT TẮT VÀ THUẬT NGỮ}

\begin{longtable}{c l p{7.6cm}}
\toprule
\textbf{STT} & \textbf{Từ viết tắt} & \textbf{Giải nghĩa}\\
\midrule
\endfirsthead
\toprule
\textbf{STT} & \textbf{Từ viết tắt} & \textbf{Giải nghĩa}\\
\midrule
\endhead
1  & API   & Application Programming Interface -- Giao diện lập trình ứng dụng\\
2  & CI/CD & Continuous Integration / Continuous Deployment -- Tích hợp và triển khai liên tục\\
3  & CSDL  & Cơ sở dữ liệu\\
4  & GitOps & Phương pháp vận hành hệ thống lấy Git làm nguồn cấu hình duy nhất\\
5  & HA    & High Availability -- Tính sẵn sàng cao\\
6  & HTTP  & Hypertext Transfer Protocol -- Giao thức truyền tải siêu văn bản\\
7  & IaC   & Infrastructure as Code -- Hạ tầng dưới dạng mã\\
8  & IdP   & Identity Provider -- Nhà cung cấp định danh\\
9  & IAM   & Identity and Access Management -- Quản lý định danh và truy cập\\
10 & JWT   & JSON Web Token\\
11 & K8s   & Kubernetes -- Nền tảng điều phối container\\
12 & MFA   & Multi-Factor Authentication -- Xác thực đa yếu tố\\
13 & OIDC  & OpenID Connect\\
14 & OAuth & Open Authorization\\
15 & PKCE  & Proof Key for Code Exchange\\
16 & RP    & Relying Party -- Bên tin cậy (ứng dụng dựa vào IdP)\\
17 & SSO   & Single Sign-On -- Đăng nhập một lần\\
18 & TLS   & Transport Layer Security -- Giao thức bảo mật tầng truyền tải\\
19 & TOTP  & Time-based One-Time Password -- Mật khẩu dùng một lần theo thời gian\\
20 & RBAC  & Role-Based Access Control -- Kiểm soát truy cập theo vai trò\\
\bottomrule
\end{longtable}

\vspace{0.5cm}
\textit{Ghi chú: Danh mục được sắp xếp theo thứ tự bảng chữ cái của từ viết tắt.
Cập nhật thêm khi bổ sung thuật ngữ mới trong quá trình viết.}
